\section{Inference Backends}
\label{sec:backends}

\subsection{Backend Architecture}

Miniforge abstracts over multiple inference backends through a unified interface, enabling platform-appropriate deployment and graceful fallback.

\subsubsection{Backend Interface}

All backends implement the common interface:

\begin{lstlisting}[language=Python, caption=Backend Interface]
class InferenceBackend(ABC):
    async def initialize(self) -> None: ...
    
    async def generate(
        self, prompt: str, max_tokens: int,
        temperature: float, top_p: float, top_k: int,
        stop: Optional[List[str]]
    ) -> str: ...
    
    async def generate_stream(
        self, prompt: str, **kwargs
    ) -> AsyncIterator[str]: ...
    
    async def get_info(self) -> Dict[str, Any]: ...
    
    async def cleanup(self) -> None: ...
\end{lstlisting}

\subsection{llama.cpp Backend}

The llama.cpp backend provides optimized CPU inference with GGUF support.

\subsubsection{Key Features}

\begin{itemize}
    \item GGUF quantization formats (Q4\_K\_M, Q5\_K\_M, etc.)
    \item KV cache compression (TurboQuant)
    \item Flash Attention
    \item Multi-threading optimized for CPU
    \item Cross-platform support
\end{itemize}

\subsubsection{Configuration Options}

Key parameters for llama.cpp backend:

\begin{lstlisting}[language=Python, caption=llama.cpp Configuration]
config = {
    "n_ctx": 194560,           # Context window
    "n_threads": 8,            # CPU threads
    "n_batch": 2048,           # Batch size
    "n_ubatch": 512,           # Micro-batch size
    "cache_type_k": "turbo3",  # Key cache quantization
    "cache_type_v": "turbo3",  # Value cache quantization
    "flash_attn": True,        # Enable Flash Attention
    "use_mmap": True,          # Memory-map weights
    "use_mlock": False,        # Disable (WSL2)
}
\end{lstlisting}

\subsubsection{Context Fallback}

When initialization fails due to memory constraints, the backend implements automatic context fallback:

\begin{enumerate}
    \item Try requested context size
    \item If OOM, try progressively smaller sizes
    \item Log warning about reduced context
    \item Continue with successful size
\end{enumerate}

This ensures the model loads even under memory pressure.

\subsection{Transformers Backend}

The Transformers backend provides native HuggingFace compatibility as a fallback.

\subsubsection{Use Cases}

\begin{itemize}
    \item When GGUF conversion is unavailable
    \item For testing and debugging
    \item When llama.cpp fails to initialize
    \item For features not yet implemented in llama.cpp
\end{itemize}

\subsubsection{Configuration}

Transformers backend uses PyTorch for inference:

\begin{lstlisting}[language=Python, caption=Transformers Configuration]
config = {
    "torch_dtype": torch.float16,
    "device_map": "auto",
    "low_cpu_mem_usage": True,
}
\end{lstlisting}

\subsubsection{Performance Comparison}

Table \ref{tab:backend_comparison} compares backend performance.

\begin{table}[h]
\centering
\caption{Backend Performance Comparison}
\label{tab:backend_comparison}
\begin{tabular}{lcc}
\toprule
\textbf{Metric} & \textbf{llama.cpp} & \textbf{Transformers} \\
\midrule
Throughput (tok/s) & 15-25 & 8-12 \\
Memory Overhead & Low & High \\
Startup Time & Fast & Slow \\
Quantization Support & Native & Limited \\
CPU Optimization & Excellent & Good \\
\bottomrule
\end{tabular}
\end{table}

\subsection{Backend Selection Strategy}

Miniforge automatically selects the appropriate backend:

\begin{enumerate}
    \item Attempt llama.cpp with requested quantization
    \item If GGUF unavailable, try auto-conversion
    \item If conversion fails, fall back to Transformers
    \item User can override with explicit \texttt{backend=} parameter
\end{enumerate}

\subsection{Backend Fallback Mechanism}

The fallback mechanism ensures robust deployment:

\begin{lstlisting}[language=Python, caption=Backend Fallback]
async def _load_model(self, quantization: str):
    try:
        # Try llama.cpp first
        self.backend_name = "llama_cpp"
        await self._try_load_llama_cpp(quantization)
    except Exception as e:
        logger.warning(f"llama.cpp failed: {e}")
        
        # Check if explicit GGUF repo
        if "unsloth" in model_id or "-GGUF" in model_id:
            raise  # Don't fallback for GGUF repos
        
        # Fallback to transformers
        logger.info("Falling back to transformers backend")
        self.backend_name = "transformers"
        await self._try_load_transformers()
\end{lstlisting}

\subsection{Unified API Benefits}

The unified backend abstraction provides:

\begin{itemize}
    \item \textbf{Portability:} Same code works across backends
    \item \textbf{Robustness:} Automatic fallback on failures
    \item \textbf{Flexibility:} Choose backend based on requirements
    \item \textbf{Testing:} Easy to compare backend performance
\end{itemize}

\subsection{Future Backend Extensions}

Potential future backends include:

\begin{itemize}
    \item \textbf{ONNX Runtime:} Cross-platform optimization
    \item \textbf{TensorRT:} If GPU becomes available
    \item \textbf{Custom Kernels:} Hand-optimized for specific hardware
\end{itemize}
